\section{Discussion and Limitations}
\label{sec:discussion}

\paragraph{Why placement, and not reconstruction, carries the gain.} The literature spends its design effort on the replacement of a skipped step, and our measurements separate that choice from the choice of where to skip. Inside one pipeline, replacing the reconstruction of a skipped step by plain reuse leaves the LPIPS of our schedule unchanged, while replacing the schedule by a threshold that keeps the same number of calls multiplies the loss by three. The two choices are not symmetric: the replacement acts on a step that has already been skipped, so it can only reduce an error that placement has already admitted, whereas placement decides which errors exist. A measured profile of the sampler tells us where those errors are cheap, and it is available before denoising starts.

\paragraph{The profile is a property of the sampler.} The order of the first eight steps is identical across prompts and the two profiles correlate at 0.94, which is what makes a fixed schedule defensible at all. The mechanism is a matter of inheritance: an error introduced while the latent is still noise is carried by every subsequent step, because each step conditions on the latent that the previous one produced, while an error introduced near the end is observed by a handful of exact evaluations that re-explain it before the decode. The same mechanism predicts the shape of the tail: since the profile decays, no two decisions near the end deserve equal spacing, which is why the gaps grow rather than stay uniform.

\paragraph{What we do not claim.} The head length is not a universal constant. It is read from a measurement that depends on the step count, the scheduler and the backbone, and we re-measure it once per setting rather than transferring it. The budgets we study lie between ten and twenty-one calls, so the regime below ten, where the head can no longer be held exact, is characterised only by the collapse of every rule we tried. The study covers two backbones, CogVideoX-2B and Wan2.1-1.3B, and two frame counts, and we do not claim a wall-clock result: the budget is counted in transformer calls, which fixes the cost of the network but leaves the decoder and the host out of the comparison. Finally, all baselines are re-run inside one implementation, so the absolute values of Table~\ref{tab:main} are not the numbers those papers report; the ordering within one pipeline is what the table supports.
